%\noindent
%\includegraphics[width=90mm]{nwb-graphics/nwbLogo.jpg}
%\\
%\\
%\noindent
%\includegraphics[width=90mm]{graphics/cishellLogo.png}

\section{Overview}

The Network Workbench (NWB) project will design, evaluate, and operate a unique 
distributed, shared resources environment for biomedical, social science, 
and physics research. Investigators are Katy B\"{o}rner, Santiago Schnell, 
Alessandro Vespignani, Stanley Wasserman, Eric Wernert (Indiana University), 
and Albert-L\'{a}szl\'{o} Barab\'{a}si (Notre Dame University). Major software 
developers are Weixia Huang, Bruce Herr, Ben Markines, Santo Fortunato (Indiana 
University), and Cesar Hidalgo (Notre Dame University). Advisory board members 
comprise Noshir Contractor (NCSA), Craig Alan Stewart (Indiana University), 
James Hendler (University of Maryland), Jason Leigh (University of Illinois at 
Chicago), Neo Martinez (Rocky Mountain Biological Lab), Michael Macy (Cornell 
University), Ulrik Brandes (University of Konstanz), Mark Gerstein (Yale 
University), Stephen North (AT\&T), and Tom Snijders (University of Groningen). 
The work is supported in part by a NSF IIS-0513650 award.
%\includegraphics[height=3mm]{nwb-graphics/nsf4c.png}

The NWB data-code-computing resources environment will provide a one-stop 
online portal for researchers, educators, and practitioners. Users of the NWB 
will have online access to major network datasets or can upload their own 
networks. They will be able to perform network analysis with the most effective 
algorithms available. In addition, they will be able to generate, run, and 
validate network models to advance their understanding of the structure and 
dynamics of particular networks. NWB will provide advanced visualization tools 
to interactively explore and understand specific networks, as well as their 
interaction with other types of networks.

A major computer science challenge is the development of an algorithm 
integration framework, called the Cyberinfrastructure Shell (CIShell), that 
supports the easy integration and dissemination of existing and new algorithms 
and can deal with the multitude of network data formats in existence today. 
Another challenge is the design and implementation of an easy to use 
menu-based, online portal interface for interactive algorithm selection, data 
manipulation, user and session management.

The NWB will be evaluated in diverse research projects and educational settings 
in biology, social and behavioral science, and physics research.

This document describes CIShell, the CIShell Reference GUI and the NWB Tool, as 
well as the NWB Community Wiki.

%\noindent
%\includegraphics[width=90mm]{nwb-graphics/PIs.png}
%\noindent
%\includegraphics[width=80mm]{nwb-graphics/developers.png}
% \noindent
% \includegraphics[width=90mm]{nwb-graphics/advisoryBoard.jpg}